\documentclass[12pt]{article}
\usepackage[brazilian]{babel}
\usepackage{csquotes}
\usepackage{glossaries}
\usepackage[
  perfil=artigo,
  artigo-colunas=uma,
  citacoes=autor-data
]{abntex3}

\makeglossaries
\newglossaryentry{artefato}{
  name={artefato digital},
  description={objeto produzido ou usado durante um processo computacional}
}
\addbibresource{referencias-exemplo.bib}
\ABNTEXsetup{
  titulo={Modelo canônico completo de artigo em uma coluna},
  subtitulo={base auditável com todos os elementos implementados}
}
\ABNTEXartigoSetup{
  titulo-outro-idioma={Complete canonical one-column article template},
  resumo={Este artigo demonstra todos os elementos públicos aplicáveis ao perfil, com autoria múltipla, dados editoriais e pós-textuais opcionais.},
  palavras-chave={artigo; normalização; auditoria; LaTeX},
  resumo-outro-idioma={This article demonstrates every public element applicable to the profile, including multiple authors, editorial data, and optional back matter.},
  palavras-chave-outro-idioma={article; standardization; audit; LaTeX},
  data-submissao={10 de janeiro de 2026},
  data-aprovacao={20 de março de 2026},
  identificacao-disponibilidade={DOI: 10.0000/exemplo.uma-coluna}
}
\ABNTEXartigoautor
  {Ana Silva}
  {Doutora em Ciência da Informação}
  {Universidade Exemplo}
  {ana@example.org}
\ABNTEXartigoautor
  {Bruno Souza}
  {Mestre em Preservação Digital}
  {Instituto Exemplo}
  {bruno@example.org}

\begin{document}
\ABNTEXartigopretextual
\ABNTEXartigotextual

\ABNTEXartigointroducao{%
  A introdução delimita o assunto, apresenta objetivos e situa o tema. Este
  modelo trata da preservação de cada \gls{artefato} e utiliza citações
  processadas pelo \texttt{biblatex-abnt} \parencite{associacao2025}.
}
\ABNTEXartigodesenvolvimento[Materiais e métodos]{%
  O método inventaria a API, compila o documento e inspeciona seus elementos.
  \begin{ABNTEXcitacaolonga}
  Trecho demonstrativo destinado à inspeção do ambiente de citação longa.
  \end{ABNTEXcitacaolonga}
  \ABNTEXnota{Nota demonstrativa do artigo em uma coluna.}

  \subsection{Resultados}
  \begin{figure}[ht]
    \centering
    \rule{0.55\linewidth}{2cm}
    \caption{Fluxo demonstrativo}
    \ABNTEXfonte{elaborada pelas autorias.}
  \end{figure}
  \begin{table}[ht]
    \centering
    \caption{Elementos inspecionados}
    \begin{tabular}{ll}
      Elemento & Resultado\\
      Bloco inicial & composto\\
      Referências & processadas
    \end{tabular}
    \ABNTEXfonte{elaborada pelas autorias.}
  \end{table}

  \subsection{Discussão}
  Os resultados são confrontados com a fundamentação adotada.
}
\ABNTEXartigoconsideracoesfinais{%
  As considerações respondem aos objetivos, registram limitações e indicam
  continuidade da investigação.
}

\ABNTEXartigopostextual
\ABNTEXbibliografia
\ABNTEXglossario
\ABNTEXappendix{Instrumento elaborado pelas autorias}
Material produzido especificamente para o estudo.
\ABNTEXannex{Documento externo associado ao artigo}
Material não elaborado pelas autorias.
\ABNTEXartigoagradecimentos{%
  Às instituições e às pessoas responsáveis pela revisão externa.
}
\end{document}
